\section*{Chapter \ref{ch:b3:deformation-mechanisms} --- Deformation Mechanisms and Shear Zones}

\begin{solution}{exo:b3:deformation-mechanisms:1}
Forsterite: $1573/2163 = 0.73$; ice: $263/273 = 0.96$. Ice is far closer to
its melting point and creeps more readily: glaciers flow at the Earth's
surface.
\end{solution}

\begin{solution}{exo:b3:deformation-mechanisms:2}
$\gamma = \tan\ang{70} = 2.75$ and $\tan\ang{80} = 5.67$. Near \ang{90} the
tangent changes very fast with the angle: one degree of error at $\psi =
\ang{84}$ changes $\gamma$ by a fifth.
\end{solution}

\begin{solution}{exo:b3:deformation-mechanisms:3}
$\sigma = (10^{3.56}/10)^{1/1.26} = 363.1^{0.794} = \qty{108}{MPa}$.
\end{solution}

\begin{solution}{exo:b3:deformation-mechanisms:4}
$E/RT = 223000/(8.314\times773.15) = 34.69$, so $\dot\varepsilon =
\num{1.1e-4}\times20^4\times e^{-34.69} = \num{1.1e-4}\times\num{1.6e5}\times
\num{8.6e-16} = \qty{1.5e-14}{s^{-1}}$.
\end{solution}

\begin{solution}{exo:b3:deformation-mechanisms:5}
$\eta = \sigma/(2\dot\varepsilon)$: $\num{65.4e6}/\num{2e-14} =
\qty{3.3e21}{Pa.s}$ at \qty{400}{\celsius} and $\num{18.0e6}/\num{2e-14} =
\qty{9.0e20}{Pa.s}$ at \qty{500}{\celsius}: the middle crust at these
temperatures is about as viscous as the mantle beneath the ice sheets.
\end{solution}

\begin{solution}{exo:b3:deformation-mechanisms:6}
$n = \log(\num{1.6e-4}/\num{e-5})/\log2 = \log16/\log2 = 4$: \omterm{def:b3:deformation-mechanisms:creep}{dislocation creep},
whose \omterm{def:b3:deformation-mechanisms:flow-law}{stress exponent} is well above 1.
\end{solution}

\begin{solution}{exo:b3:deformation-mechanisms:7}
$\ln10 = (E/R)(1/1000 - 1/1100)$, so $E = 8.314\times2.3026/\num{9.091e-5} =
\qty{211}{kJ.mol^{-1}}$.
\end{solution}

\begin{solution}{exo:b3:deformation-mechanisms:8}
Nabarro--Herring ($d^{-2}$): 100 times faster; Coble ($d^{-3}$): 1000 times.
Along the boundaries the diffusing atoms use a path whose cross-section, $\delta
d$, shrinks with the grain, on top of the shorter distance and the smaller
volume to move.
\end{solution}

\begin{solution}{exo:b3:deformation-mechanisms:9}
$\sigma_b^{2.5} = \num{3.59e-6}\,d^{-3}/\num{6.54e-17}$: for $d = 1000$,
$\sigma_b = 54.9^{0.4} = \qty{5.0}{MPa}$; for $d = 100$, $\sigma_b =
\num{54900}^{0.4} = \qty{79}{MPa}$. The slope is $-3/2.5 = -1.2$: a tenfold
finer grain raises the boundary by $10^{1.2} = 16$.
\end{solution}

\begin{solution}{exo:b3:deformation-mechanisms:10}
In both paths the driving difference of concentration is the same and the
distance is $d$; the volume to move per unit strain is $d^3$. The lattice
offers a cross-section $\propto d^2$, so the rate goes as $d^2/(d\,d^3) =
d^{-2}$; the boundaries offer $\propto\delta d$, so $\delta d/(d\,d^3) = \delta
d^{-3}$. The ratio of boundary to lattice rates is $(B_B\delta D_B)/(B_LD_Ld)$,
larger than 1 when $d < (B_B/B_L)\,\delta D_B/D_L$.
\end{solution}

\begin{solution}{exo:b3:deformation-mechanisms:11}
$\eta \propto \dot\varepsilon^{(1-n)/n} = \dot\varepsilon^{-2.5/3.5}$: ten times
faster divides the viscosity by $10^{0.714} = 5.2$. A band that happens to
deform faster becomes weaker, so it takes more of the deformation: the strain
gathers into it.
\end{solution}

\begin{solution}{exo:b3:deformation-mechanisms:12}
Before: $\dot\varepsilon_{\mathrm{dis}} = \num{6.54e-17}\times10^{3.5} =
\qty{2.07e-13}{s^{-1}}$ and $\dot\varepsilon_{\mathrm{dif}} = \num{3.59e-6}
\times10/3000^3 = \qty{1.3e-15}{s^{-1}}$: \qty{2.08e-13}{s^{-1}}, \omterm{def:b3:deformation-mechanisms:creep}{dislocation
creep}. After: $\dot\varepsilon_{\mathrm{dif}} = \num{3.59e-6}\times10/30^3 =
\qty{1.33e-9}{s^{-1}}$, \omterm{def:b3:deformation-mechanisms:creep}{diffusion creep}, \num{6400} times faster at the same
stress: recrystallisation has softened the zone enormously.
\end{solution}

\begin{solution}{pb:b3:deformation-mechanisms:1}
\textbf{1.} The dyke turns from \ang{90} to \ang{6} from the zone: $\psi =
\ang{84}$ and $\gamma = \tan\ang{84} = 9.51$.

\textbf{2.} $\gamma w = 9.51\times400 = \qty{3804}{m}$, about \qty{3.8}{km}.

\textbf{3.} $S_1 = \sqrt{1 + 9.51^2/4} + 9.51/2 = 4.859 + 4.755 = 9.614$,
$S_2 = 1/S_1 = 0.104$, and $S_1/S_2 = 92.4$.

\textbf{4.} $\tan2\theta' = 2/9.51 = 0.210$, $\theta' = \ang{5.9}$: the
foliation, close to the long axis, turns ever closer to the zone as $\gamma$
grows.

\textbf{5.} Followed northwards, the dyke is carried east inside the zone and
leaves it further east: the northern wall moved east relative to the southern.
Seen from the southern wall, the far side moved to the right: the zone is
dextral.

\textbf{6.} In the plane containing the stretching lineation, horizontal here
for a strike-slip zone, and perpendicular to the vertical foliation: a
horizontal section. Its S foliation leans anticlockwise from the zone, towards
east-north-east, like the long axis of the strain ellipse.

\textbf{7.} $\sigma = (10^{3.56}/20)^{1/1.26} = 181.5^{0.794} = \qty{62}{MPa}$.

\textbf{8.} $242.1^{0.794} = \qty{78}{MPa}$ for \qty{15}{\micro m} and
$145.2^{0.794} = \qty{52}{MPa}$ for \qty{25}{\micro m}.

\textbf{9.} Yes: the experiments spanned 34 to \qty{268}{MPa}.

\textbf{10.} The zone deformed at \qty{450}{\celsius}, far below the 800 to
\qty{1100}{\celsius} of the experiments; a relation that does not depend on
the temperature can be carried down to it, though it is still an
extrapolation.

\textbf{11.} $D = 3631/100^{1.26} = 3631/331 = \qty{11}{\micro m}$.

\textbf{12.} $E/RT = 223000/(8.314\times723.15) = 37.09$ and $e^{-37.09} =
\num{7.79e-17}$.

\textbf{13.} $\dot\varepsilon = \num{1.1e-4}\times62^4\times\num{7.79e-17} =
\num{1.1e-4}\times\num{1.478e7}\times\num{7.79e-17} = \qty{1.27e-13}{s^{-1}}$.

\textbf{14.} $\dot\gamma = 2\dot\varepsilon = \qty{2.53e-13}{s^{-1}}$.

\textbf{15.} $\eta = \num{62e6}/(2\times\num{1.27e-13}) = \qty{2.4e20}{Pa.s}$,
close to the \qty{3e20}{Pa.s} of the upper mantle beneath Lake Ardo: the zone
was as weak as the mantle.

\textbf{16.} $e^{-E/RT}$ changes by the factor 0.064 at \qty{400}{\celsius}
(sixteen times slower) and 11.0 at \qty{500}{\celsius}.

\textbf{17.} $(78/62)^4 = 2.5$ times faster.

\textbf{18.} $t = 9.51/\num{2.53e-13} = \qty{3.75e13}{s}$.

\textbf{19.} $\num{3.75e13}/\num{3.156e10} = \qty{1190}{kyr}$.

\textbf{20.} $\qty{3804}{m}/\qty{1.19e6}{yr} = \qty{3.2}{mm/yr}$.

\textbf{21.} Ten or more times slower than the main plate boundaries: the zone
was a minor strand, or one of several that shared a boundary's motion.

\textbf{22.} $\dot\gamma = 0.02/(\num{3.156e7}\times400) = \qty{1.58e-12}{s^{-1}}$,
6.2 times faster; the stress rises by $6.2^{1/4} = 1.58$, to \qty{98}{MPa}, and
the grains would be $3631/98^{1.26} = \qty{11}{\micro m}$.

\textbf{23.} The time goes as $\sigma^{-4}$: the grain-size range alone gives
$(62/78)^4 = 0.40$ to $(62/52)^4 = 2.0$ times the estimate, and the
uncertainty of the exponent $1.26 \pm 0.13$ and of the \omterm{def:b3:deformation-mechanisms:flow-law}{flow law}'s constants
widens it further: the result is an order of magnitude, not a date.

\textbf{24.} The Varden \omterm{def:b3:deformation-mechanisms:shear-zone}{shear zone} took about
\textbf{\qty[text-series-to-math=true]{1190}{kyr}} to accumulate its strain.
\end{solution}
